\providecommand\ExamFontPath{../../../字体/}
\documentclass[fontset=exam]{exam-zh}
\usepackage{longtable,booktabs,array,calc,multirow,fvextra}
\usepackage[normalem]{ulem}
\providecommand{\hl}[1]{#1}
\providecommand{\ul}[1]{\uline{#1}}
\DefineVerbatimEnvironment{verbatim}{Verbatim}{breaklines=true,fontsize=\small}
\providecommand{\tightlist}{\setlength{\itemsep}{0pt}\setlength{\parskip}{0pt}}
\newcommand{\pandocbounded}[1]{#1}
\setlength{\emergencystretch}{3em}
\setlength{\tabcolsep}{2pt}
\examsetup{page/size=a4paper}
\title{期末考试 B 卷}
\subject[4em]{离散数学}
\begin{document}
\maketitle
\fancyfoot[C]{第 \thepage 页（共 \pageref*{LastPage} 页）}
\pdfbookmark[1]{一、证明题（10分）}{section18}

\par\medskip{\heiti

一、证明题（10 分）

\par}\smallskip

1)((P\(\lor\)Q)\(\land\)\(\neg\)(\(\neg\)P\(\land\)(\(\neg\)Q\(\lor\)\(\neg\)R)))\(\lor\)(\(\neg\)P\(\land\)\(\neg\)Q)\(\lor\)(\(\neg\)P\(\land\)\(\neg\)R)\(\Leftrightarrow\)T

2)\(\forall\)x(P(x)\(\to\)Q(x))\(\land\)\(\forall\)xP(x)\(\Leftrightarrow\)\(\forall\)x(P(x)\(\land\)Q(x))

\pdfbookmark[1]{二、求命题公式(¬P→Q)→(P∨¬Q) 的主析取范式和主合取范式（10分）}{section21}

二、求命题公式(\(\neg\)P\(\to\)Q)\(\to\)(P\(\lor\)\(\neg\)Q) 的主析取范式和主合取范式（10 分）

\pdfbookmark[1]{三、推理证明题（10分）}{section22}

三、推理证明题（10 分）

1)(P\(\to\)(Q\(\to\)S))\(\land\)(\(\neg\)R\(\lor\)P)\(\land\)Q\(\Rightarrow\)R\(\to\)S

2) \(\forall\)x(P(x)\(\lor\)Q(x))，\(\forall\)x\(\neg\)P(x)\(\Rightarrow\)\(\exists\)x Q(x)

\pdfbookmark[1]{四、例5在边长为1的正方形内任意放置九个点，证明其中必存在三个点，使得由它们组成的三角形（可能是退化的）面积不超过1/8}{section25}

四、例 5 在边长为 1 的正方形内任意放置九个点，证明其中必存在三个点，使得由它们组成的三角形（可能是退化的）面积不超过 1/8（10 分）。

\pdfbookmark[1]{五、已知A、B、C是三个集合，证明A∩(B∪C)=(A∩B)∪(A∩C) （10分）}{section26}

五、已知 A、B、C 是三个集合，证明 A\(\cap\)(B\(\cup\)C)=(A\(\cap\)B)\(\cup\)(A\(\cap\)C) （10 分）

\pdfbookmark[1]{六、π=A1，A2，…，An是集合A的一个划分，定义R=<a，b>|a、b∈Ai，I=1，2，…，n，则R是A上的等价关}{section27}

六、π=\{A\textsubscript{1}，A\textsubscript{2}，\ldots，A\textsubscript{n}\}是集合 A 的一个划分，定义 R=\{\textless a，b\textgreater\textbar a、b\(\in\)A\textsubscript{i}，I=1，2，\ldots，n\}，则 R 是 A 上的等价关系（15 分）。

\pdfbookmark[1]{七、若f:A→B是双射，则f-1:B→A是双射（15分）。}{section28}

七、若 f:A\(\to\)B 是双射，则 f\textsuperscript{-1}:B\(\to\)A 是双射（15 分）。

\pdfbookmark[1]{八、设<G，*>是群，<A，*>和<B，*>是<G，*>的子群，证明：若A∪B＝G，则A＝G或B＝G（10分）。}{section29}

八、设\textless{}\emph{G}，*\textgreater 是群，\textless{}\emph{A}，*\textgreater 和\textless{}\emph{B}，*\textgreater 是\textless{}\emph{G}，*\textgreater 的子群，证明：若\emph{A}\(\cup\)\emph{B}＝\emph{G}，则\emph{A}＝\emph{G}或\emph{B}＝\emph{G}（10 分）。

\pdfbookmark[1]{九、若无向图G是不连通的，证明G的补图overlineG是连通的（10分）。}{section30}

九、若无向图\emph{G}是不连通的，证明\emph{G}的补图\(\overline{G}\)是连通的（10 分）。

\end{document}
